%% bare_conf.tex
%% V1.4
%% 2012/12/27
%% by Michael Shell
%% See:
%% http://www.michaelshell.org/
%% for current contact information.
%%
%% This is a skeleton file demonstrating the use of IEEEtran.cls
%% (requires IEEEtran.cls version 1.8 or later) with an IEEE conference paper.
%%
%% Support sites:
%% http://www.michaelshell.org/tex/ieeetran/
%% http://www.ctan.org/tex-archive/macros/latex/contrib/IEEEtran/
%% and
%% http://www.ieee.org/

%%*************************************************************************
%% Legal Notice:
%% This code is offered as-is without any warranty either expressed or
%% implied; without even the implied warranty of MERCHANTABILITY or
%% FITNESS FOR A PARTICULAR PURPOSE! 
%% User assumes all risk.
%% In no event shall IEEE or any contributor to this code be liable for
%% any damages or losses, including, but not limited to, incidental,
%% consequential, or any other damages, resulting from the use or misuse
%% of any information contained here.
%%
%% All comments are the opinions of their respective authors and are not
%% necessarily endorsed by the IEEE.
%%
%% This work is distributed under the LaTeX Project Public License (LPPL)
%% ( http://www.latex-project.org/ ) version 1.3, and may be freely used,
%% distributed and modified. A copy of the LPPL, version 1.3, is included
%% in the base LaTeX documentation of all distributions of LaTeX released
%% 2003/12/01 or later.
%% Retain all contribution notices and credits.
%% ** Modified files should be clearly indicated as such, including  **
%% ** renaming them and changing author support contact information. **
%%
%% File list of work: IEEEtran.cls, IEEEtran_HOWTO.pdf, bare_adv.tex,
%%                    bare_conf.tex, bare_jrnl.tex, bare_jrnl_compsoc.tex,
%%                    bare_jrnl_transmag.tex
%%*************************************************************************

% *** Authors should verify (and, if needed, correct) their LaTeX system  ***
% *** with the testflow diagnostic prior to trusting their LaTeX platform ***
% *** with production work. IEEE's font choices can trigger bugs that do  ***
% *** not appear when using other class files.                            ***
% The testflow support page is at:
% http://www.michaelshell.org/tex/testflow/



% Note that the a4paper option is mainly intended so that authors in
% countries using A4 can easily print to A4 and see how their papers will
% look in print - the typesetting of the document will not typically be
% affected with changes in paper size (but the bottom and side margins will).
% Use the testflow package mentioned above to verify correct handling of
% both paper sizes by the user's LaTeX system.
%
% Also note that the "draftcls" or "draftclsnofoot", not "draft", option
% should be used if it is desired that the figures are to be displayed in
% draft mode.
%
\documentclass[conference]{IEEEtran}
% Add the compsoc option for Computer Society conferences.
%
% If IEEEtran.cls has not been installed into the LaTeX system files,
% manually specify the path to it like:
% \documentclass[conference]{../sty/IEEEtran}



% Some very useful LaTeX packages include:
% (uncomment the ones you want to load)


% For directly writing german umlauts uncomment the appropriate line for
% your operating system:
% Windows:
% \usepackage[ansinew]{inputenc}
% Linux:
\usepackage[latin1]{inputenc}
% Mac
% \usepackage[applemac]{inputenc}
% If none of the above lines work you can also try the following:
% \usepackage[utf8]{inputenc}



% *** MISC UTILITY PACKAGES ***
%
%\usepackage{ifpdf}
% Heiko Oberdiek's ifpdf.sty is very useful if you need conditional
% compilation based on whether the output is pdf or dvi.
% usage:
% \ifpdf
%   % pdf code
% \else
%   % dvi code
% \fi
% The latest version of ifpdf.sty can be obtained from:
% http://www.ctan.org/tex-archive/macros/latex/contrib/oberdiek/
% Also, note that IEEEtran.cls V1.7 and later provides a builtin
% \ifCLASSINFOpdf conditional that works the same way.
% When switching from latex to pdflatex and vice-versa, the compiler may
% have to be run twice to clear warning/error messages.





% *** CITATION PACKAGES ***
%
%\usepackage{cite}
% cite.sty was written by Donald Arseneau
% V1.6 and later of IEEEtran pre-defines the format of the cite.sty package
% \cite{} output to follow that of IEEE. Loading the cite package will
% result in citation numbers being automatically sorted and properly
% "compressed/ranged". e.g., [1], [9], [2], [7], [5], [6] without using
% cite.sty will become [1], [2], [5]--[7], [9] using cite.sty. cite.sty's
% \cite will automatically add leading space, if needed. Use cite.sty's
% noadjust option (cite.sty V3.8 and later) if you want to turn this off
% such as if a citation ever needs to be enclosed in parenthesis.
% cite.sty is already installed on most LaTeX systems. Be sure and use
% version 4.0 (2003-05-27) and later if using hyperref.sty. cite.sty does
% not currently provide for hyperlinked citations.
% The latest version can be obtained at:
% http://www.ctan.org/tex-archive/macros/latex/contrib/cite/
% The documentation is contained in the cite.sty file itself.





% *** GRAPHICS RELATED PACKAGES ***
%
\ifCLASSINFOpdf
  \usepackage[pdftex]{graphicx}
  % declare the path(s) where your graphic files are
  \graphicspath{{../notebooks/figures/}}
  % and their extensions so you won't have to specify these with
  % every instance of \includegraphics
  % \DeclareGraphicsExtensions{.pdf,.jpeg,.png}
\else
  % or other class option (dvipsone, dvipdf, if not using dvips). graphicx
  % will default to the driver specified in the system graphics.cfg if no
  % driver is specified.
  \usepackage[dvips]{graphicx}
  % declare the path(s) where your graphic files are
  \graphicspath{{../notebooks/figures/}}
  % and their extensions so you won't have to specify these with
  % every instance of \includegraphics
  % \DeclareGraphicsExtensions{.eps}
\fi
% graphicx was written by David Carlisle and Sebastian Rahtz. It is
% required if you want graphics, photos, etc. graphicx.sty is already
% installed on most LaTeX systems. The latest version and documentation
% can be obtained at: 
% http://www.ctan.org/tex-archive/macros/latex/required/graphics/
% Another good source of documentation is "Using Imported Graphics in
% LaTeX2e" by Keith Reckdahl which can be found at:
% http://www.ctan.org/tex-archive/info/epslatex/
%
% latex, and pdflatex in dvi mode, support graphics in encapsulated
% postscript (.eps) format. pdflatex in pdf mode supports graphics
% in .pdf, .jpeg, .png and .mps (metapost) formats. Users should ensure
% that all non-photo figures use a vector format (.eps, .pdf, .mps) and
% not a bitmapped formats (.jpeg, .png). IEEE frowns on bitmapped formats
% which can result in "jaggedy"/blurry rendering of lines and letters as
% well as large increases in file sizes.
%
% You can find documentation about the pdfTeX application at:
% http://www.tug.org/applications/pdftex





% *** MATH PACKAGES ***
%
%\usepackage[cmex10]{amsmath}
% A popular package from the American Mathematical Society that provides
% many useful and powerful commands for dealing with mathematics. If using
% it, be sure to load this package with the cmex10 option to ensure that
% only type 1 fonts will utilized at all point sizes. Without this option,
% it is possible that some math symbols, particularly those within
% footnotes, will be rendered in bitmap form which will result in a
% document that can not be IEEE Xplore compliant!
%
% Also, note that the amsmath package sets \interdisplaylinepenalty to 10000
% thus preventing page breaks from occurring within multiline equations. Use:
%\interdisplaylinepenalty=2500
% after loading amsmath to restore such page breaks as IEEEtran.cls normally
% does. amsmath.sty is already installed on most LaTeX systems. The latest
% version and documentation can be obtained at:
% http://www.ctan.org/tex-archive/macros/latex/required/amslatex/math/





% *** SPECIALIZED LIST PACKAGES ***
%
%\usepackage{algorithmic}
% algorithmic.sty was written by Peter Williams and Rogerio Brito.
% This package provides an algorithmic environment fo describing algorithms.
% You can use the algorithmic environment in-text or within a figure
% environment to provide for a floating algorithm. Do NOT use the algorithm
% floating environment provided by algorithm.sty (by the same authors) or
% algorithm2e.sty (by Christophe Fiorio) as IEEE does not use dedicated
% algorithm float types and packages that provide these will not provide
% correct IEEE style captions. The latest version and documentation of
% algorithmic.sty can be obtained at:
% http://www.ctan.org/tex-archive/macros/latex/contrib/algorithms/
% There is also a support site at:
% http://algorithms.berlios.de/index.html
% Also of interest may be the (relatively newer and more customizable)
% algorithmicx.sty package by Szasz Janos:
% http://www.ctan.org/tex-archive/macros/latex/contrib/algorithmicx/




% *** ALIGNMENT PACKAGES ***
%
\usepackage{array}
% Frank Mittelbach's and David Carlisle's array.sty patches and improves
% the standard LaTeX2e array and tabular environments to provide better
% appearance and additional user controls. As the default LaTeX2e table
% generation code is lacking to the point of almost being broken with
% respect to the quality of the end results, all users are strongly
% advised to use an enhanced (at the very least that provided by array.sty)
% set of table tools. array.sty is already installed on most systems. The
% latest version and documentation can be obtained at:
% http://www.ctan.org/tex-archive/macros/latex/required/tools/


% IEEEtran contains the IEEEeqnarray family of commands that can be used to
% generate multiline equations as well as matrices, tables, etc., of high
% quality.




% *** SUBFIGURE PACKAGES ***
%\ifCLASSOPTIONcompsoc
%  \usepackage[caption=false,font=normalsize,labelfont=sf,textfont=sf]{subfig}
%\else
%  \usepackage[caption=false,font=footnotesize]{subfig}
%\fi
% subfig.sty, written by Steven Douglas Cochran, is the modern replacement
% for subfigure.sty, the latter of which is no longer maintained and is
% incompatible with some LaTeX packages including fixltx2e. However,
% subfig.sty requires and automatically loads Axel Sommerfeldt's caption.sty
% which will override IEEEtran.cls' handling of captions and this will result
% in non-IEEE style figure/table captions. To prevent this problem, be sure
% and invoke subfig.sty's "caption=false" package option (available since
% subfig.sty version 1.3, 2005/06/28) as this is will preserve IEEEtran.cls
% handling of captions.
% Note that the Computer Society format requires a larger sans serif font
% than the serif footnote size font used in traditional IEEE formatting
% and thus the need to invoke different subfig.sty package options depending
% on whether compsoc mode has been enabled.
%
% The latest version and documentation of subfig.sty can be obtained at:
% http://www.ctan.org/tex-archive/macros/latex/contrib/subfig/




% *** FLOAT PACKAGES ***
%
%\usepackage{fixltx2e}
% fixltx2e, the successor to the earlier fix2col.sty, was written by
% Frank Mittelbach and David Carlisle. This package corrects a few problems
% in the LaTeX2e kernel, the most notable of which is that in current
% LaTeX2e releases, the ordering of single and double column floats is not
% guaranteed to be preserved. Thus, an unpatched LaTeX2e can allow a
% single column figure to be placed prior to an earlier double column
% figure. The latest version and documentation can be found at:
% http://www.ctan.org/tex-archive/macros/latex/base/


%\usepackage{stfloats}
% stfloats.sty was written by Sigitas Tolusis. This package gives LaTeX2e
% the ability to do double column floats at the bottom of the page as well
% as the top. (e.g., "\begin{figure*}[!b]" is not normally possible in
% LaTeX2e). It also provides a command:
%\fnbelowfloat
% to enable the placement of footnotes below bottom floats (the standard
% LaTeX2e kernel puts them above bottom floats). This is an invasive package
% which rewrites many portions of the LaTeX2e float routines. It may not work
% with other packages that modify the LaTeX2e float routines. The latest
% version and documentation can be obtained at:
% http://www.ctan.org/tex-archive/macros/latex/contrib/sttools/
% Do not use the stfloats baselinefloat ability as IEEE does not allow
% \baselineskip to stretch. Authors submitting work to the IEEE should note
% that IEEE rarely uses double column equations and that authors should try
% to avoid such use. Do not be tempted to use the cuted.sty or midfloat.sty
% packages (also by Sigitas Tolusis) as IEEE does not format its papers in
% such ways.
% Do not attempt to use stfloats with fixltx2e as they are incompatible.
% Instead, use Morten Hogholm'a dblfloatfix which combines the features
% of both fixltx2e and stfloats:
%
% \usepackage{dblfloatfix}
% The latest version can be found at:
% http://www.ctan.org/tex-archive/macros/latex/contrib/dblfloatfix/




% *** PDF, URL AND HYPERLINK PACKAGES ***
%
%\usepackage{url}
% url.sty was written by Donald Arseneau. It provides better support for
% handling and breaking URLs. url.sty is already installed on most LaTeX
% systems. The latest version and documentation can be obtained at:
% http://www.ctan.org/tex-archive/macros/latex/contrib/url/
% Basically, \url{my_url_here}.




% *** Do not adjust lengths that control margins, column widths, etc. ***
% *** Do not use packages that alter fonts (such as pslatex).         ***
% There should be no need to do such things with IEEEtran.cls V1.6 and later.
% (Unless specifically asked to do so by the journal or conference you plan
% to submit to, of course. )


% add custom packages

% lipsum
\usepackage{xcolor}
\usepackage{lipsum}
\setlipsum{%
    par-before = \begingroup\color{lightgray},
    par-after = \endgroup,
    sentence-before = \begingroup\color{lightgray},
    sentence-after = \endgroup
}

\definecolor{tumblue}{rgb}{0, 0.4, 0.74}

\newcommand{\comment}[1]{{\color{tumblue} #1}}



% correct bad hyphenation here
\hyphenation{op-tical net-works semi-conduc-tor}

\usepackage{amsfonts}
\usepackage{amssymb}
\usepackage{amsmath}
\usepackage{amsthm}
\newtheorem{theorem}{Theorem}
\newtheorem{proposition}{Proposition}
\newtheorem{example}{Example}
\newtheorem{definition}{Definition}
\newtheorem{lemma}{Lemma}

\usepackage[colorlinks=true,linkcolor=blue!70!black,
  citecolor=green!50!black,urlcolor=blue!70!black]{hyperref}
\usepackage{cleveref}

\crefname{section}{Sec.}{Sec.}
\crefname{subsection}{Sec.}{Sec.}
\crefname{figure}{Fig.}{Fig.}
\crefname{algorithm}{Alg.}{Alg.}
\crefname{table}{Tab.}{Tab.}
\crefname{example}{Example}{Example}
\crefname{definition}{Def.}{Def.}
\crefname{proposition}{Prop.}{Prop.}
\crefname{theorem}{Thm.}{Thm.}
\crefname{lemma}{Lemma}{Lemmas}

\Crefname{section}{Sec.}{Sec.}
\Crefname{subsection}{Sec.}{Sec.}
\Crefname{figure}{Fig.}{Fig.}
\Crefname{algorithm}{Alg.}{Alg.}
\Crefname{table}{Tab.}{Tab.}
\crefname{example}{Example}{Example}
\Crefname{definition}{Def.}{Def.}
\Crefname{proposition}{Prop.}{Prop.}
\Crefname{theorem}{Thm.}{Thm.}
\Crefname{lemma}{Lemma}{Lemmas}

\usepackage{booktabs} % professional tables


% tikz
% you can export matlab figures directly to latex/tikz using the following setup.
% Ask your supervisor for access to our matlab2tikz repository if you want to use it
% Important: You have to add --shell-escape to your latex call for tikz externalization to work properly
\usepackage{tikzscale}
\usepackage{pgfplots}
\usepackage{mathtools}
\pgfplotsset{compat=1.18}
\usepgfplotslibrary{external}
\tikzexternalize[prefix=./figures/externalize/]
\newcommand{\includetikz}[1]{%
    \tikzsetnextfilename{#1}%
    \input{#1.tikz}%
}
% explicitly set line width to avoid different widths in different pdf viewers
\tikzset{every picture/.style={line width=0.5pt}}
% groupplots
\usepgfplotslibrary{groupplots}
\usetikzlibrary{pgfplots.groupplots}
\usetikzlibrary{matrix}


% for operators and sets
\newcommand{\operator}[2]{\text{\normalfont\small\texttt{#1}}\left(#2\right)}
\newcommand{\opDiag}[1]{\operator{diag}{#1}}

\usepackage{xifthen}
\newcommand{\negphantom}[1]{\settowidth{\dimen0}{#1}\hspace*{-\dimen0}} % used to position punctuation neatly
\newcommand{\defSet}[3][2]{\left<#3\right>_{#2}\ifthenelse{\isempty{#1}}{}{\negphantom{#2}#1\phantom{\scriptstyle #2}}}
\newcommand{\defSetAligned}[3]{\begin{split}#2 = \left<#3\right>_{#1}\end{split}}

% Additional packages and commands for the BNN report
\usepackage{bm}
\usepackage{float}
\usepackage{enumitem}
\usepackage{listings}
\usepackage{caption}
\usepackage{subcaption}
\usepackage{rotating}
\usepackage[numbers,sort&compress]{natbib}
\usepackage{fancyhdr}
\usepackage{parskip}

\definecolor{codebg}{RGB}{248,248,248}
\lstdefinestyle{pythonstyle}{
  backgroundcolor=\color{codebg},basicstyle=\ttfamily\small,
  keywordstyle=\bfseries,frame=single,breaklines=true,
  language=Python,showstringspaces=false,
  numbers=left,numberstyle=\tiny\color{gray},xleftmargin=2em}
\lstset{style=pythonstyle}

\newcommand{\btheta}{\bm{\theta}}
\newcommand{\cD}{\mathcal{D}}
\newcommand{\cN}{\mathcal{N}}
\newcommand{\KL}{\mathrm{KL}}
\newcommand{\ELBO}{\mathcal{L}_{\text{ELBO}}}
\newcommand{\MAP}{\text{MAP}}
\newcolumntype{C}[1]{>{\centering\arraybackslash}p{#1}}
\newcolumntype{L}[1]{>{\raggedright\arraybackslash}p{#1}}

\raggedbottom

\begin{document}

% Add the seminar's cover page
\begin{figure*}[!h]

  \includegraphics{./figures/IN.pdf} \hfill \includegraphics{./figures/tumlogo.pdf}
 
  \vspace*{1cm}
  {\large \textsf{TUM School of Computation, Information and Technology}}\\
  {\large \textsf{Professorship of Cyber Physical Systems}}\\
   

  \vspace*{5cm}
%
%
% TITEL DER ARBEIT
%
%
  {\color{tumblue} \Huge \bf \textsf{Bayesian Neural Networks:
Implementations using Python Libraries}}\\  % HIER EINSETZEN!

  \vspace*{1cm}
%
%
% NAME DES STUDENTEN (auf Titelblatt)
%
% 
  {\Large \bf \textsf{Gabriel Gerardo Carpio Saez de Nanclares}}\\   % HIER EINSETZEN!
 
  \vspace*{8cm}
  {\Large \textsf{Practical Course: \emph{Formal Methods for AI-Enabled Cyber-Physical Systems} - SS 2026}}\\
 
  \vspace*{1cm} 
  \begin{tabular}{ll}
%
%
% NAME DES BETREUERS
%
%
    {\Large \bf \textsf{Advisor:}} &
    {\Large \textsf{Laura Lützow \& Victor Schulte}}\\                  % HIER EINSETZEN!
    \\

    {\Large \bf \textsf{Supervisor:}} &
    {\Large \textsf{Prof.~Dr.-Ing. Matthias Althoff}}\\
    \\

%
%
% ABGABETERMIN
%
%
    {\Large \bf \textsf{Submission:}} &
    {\Large \textsf{01. January 2025}}

  \end{tabular}
  
\end{figure*}


%
% paper title
% can use linebreaks \\ within to get better formatting as desired
% Do not put math or special symbols in the title.
\title{Bayesian Neural Networks:\\[4pt] Implementations using Python Libraries}


% author names and affiliations
% use a multiple column layout for up to three different
% affiliations
\author{\IEEEauthorblockN{Gabriel Gerardo Carpio Saez de Nanclares}
	\IEEEauthorblockA{Technische Universit\"at M\"unchen\\
		Email: go82duw@mytum.de}}

% conference papers do not typically use \thanks and this command
% is locked out in conference mode. If really needed, such as for
% the acknowledgment of grants, issue a \IEEEoverridecommandlockouts
% after \documentclass

% use for special paper notices
%\IEEEspecialpapernotice{(Invited Paper)}



% make the title area
\maketitle

% As a general rule, do not put math, special symbols or citations
% in the abstract
\begin{abstract}
	abstract

\end{abstract}

% no keywords




% For peer review papers, you can put extra information on the cover
% page as needed:
% \ifCLASSOPTIONpeerreview
% \begin{center} \bfseries EDICS Category: 3-BBND \end{center}
% \fi
%
% For peerreview papers, this IEEEtran command inserts a page break and
% creates the second title. It will be ignored for other modes.
\IEEEpeerreviewmaketitle


\section{Introduction}
\label{sec:intro}

Deep Learning finds use practically everywhere. Standard Neural Networks are
trained on given datasets in a frequentist maner, and their outputs are not
guaranteed to be true. This and the fact that the parameters set during training
are deterministic leads to over-confidence. Bayesian Deep Learning (BDL) addresses this
by replacing a deterministic weight estimate with a posterior distribution
$p(\btheta \mid \cD)$, propagating weight uncertainty into calibrated predictive distributions.
This report compares two PyTorch libraries implementing complementary Bayesian Neural Networks (BNNs) paradigms:
\begin{itemize}[leftmargin=*]
	\item \texttt{bayesian-torch}~\citep{intel2021bayesiantorch}, a
	      Bayesian by design library from Intel Labs implementing mean field
	      Variational Inference via Flipout~\citep{wen2018flipout} and
	      Bayes by Backprop~\citep{blundell2015weight}.
	\item \texttt{laplace-torch}~\citep{daxberger2021laplace,aleximmer2024laplace},
	      a library that allows applying the Laplace Approximation to any
	      pre trained model without modifying its training loop.
\end{itemize}

Experiments span three settings: two toy datasets, one is a regression
experiment, another a classification one, as well as a repli cation of the PovertyMap experiment done by Daxberger et al.~\citep{daxberger2021laplace}.
Full metric tables and figure galleries for completed experiments are in the Appendix~\ref{sec:app}.

\section{Theoretical Background}
\label{sec:theory}

\textbf{Neural Networks (NNs)~\citep{goan2020bayesian}.} A NN is a function approximator composed of layers of
interconnected neurons. There are multiple architectures, here we discuss the basis, the Multi-Layer
Perceptron (MLP), consisting of an input layer, one or more hidden layers, and
an output layer. The network is parameterized by the set $\btheta = \{\mathbf{W}^{(l)}, \mathbf{b}^{(l)}\}_{l=1}^{L}$,
where $\mathbf{W}^{(l)}$ and $\mathbf{b}^{(l)}$ are the weight matrix and bias vector for layer $l$, respectively.
Each neuron $i$ in layer $l$ computes a weighted sum of the activations from the previous layer, denoted
$\mathbf{h}^{(l-1)}$, passed through a nonlinear activation function $\sigma$:\begin{equation}h_i^{(l)} =
	\sigma!\left(\sum_j w_{ij}^{(l)} h_j^{(l-1)} + b_i^{(l)}\right),\end{equation}where $w_{ij}^{(l)}$ denotes the element
in the $i$-th row and $j$-th column of $\mathbf{W}^{(l)}$.

Training minimizes a loss function $\mathcal{L}(\btheta)$ (e.g., mean squared
error, cross-entropy) via backpropagation and
gradient-based optimization (e.g., SGD, Adam).

The libraries discussed in this report, allow implementations in the various architectures, like
Convolutional Neural Networks (CNNs), Recurrent Neural Networks (RNNs) and Transformers.

\textbf{Bayesian Neural Networks (BNNs)~\citep{goan2020bayesian,murphy2023pml}.} In Bayesian Deep Learning, Given $\cD = \{(\bm{x}_n, y_n)\}_{n=1}^{N}$ and network $f_{\btheta}$, the posterior is $p(\btheta \mid \cD) = p(\cD|\btheta)\,p(\btheta)/p(\cD)$. The evidence is intractable for modern DNNs. With a Gaussian prior
$p(\btheta)=\cN(\bm{0},\tau^{-1}\mathbf{I})$, the MAP estimate $\btheta_{\MAP} = \arg\max_{\btheta}[\log p(\cD|\btheta)+\log p(\btheta)]$
corresponds to $L_2$ regularised training with weight decay
$\lambda=\tau$~\citep{murphy2023pml}.

\subsection{Stochastic Variational Inference}

VI approximates the posterior with a tractable family $q_{\bm\phi}(\btheta)$
by maximising the Evidence Lower Bound~\citep{blundell2015weight,goan2020bayesian}:
\begin{equation}
	\ELBO(\bm\phi)
	= \mathbb{E}_{q}\!\bigl[\log p(\cD|\btheta)\bigr]
	- \KL\bigl(q_{\bm\phi}\,\|\,p(\btheta)\bigr).
	\label{eq:elbo}
\end{equation}
A mean field Gaussian $q_{\bm\phi}(w_i) = \cN(\mu_i,\sigma_i^2)$ with
$\sigma_i = \log(1+\exp(\rho_i))>0$ keeps the KL in~\eqref{eq:elbo} in
closed form, but doubles the parameter count.

\subsubsection{Bayes by Backprop}

\citet{blundell2015weight} reparameterise each weight as
\begin{equation}
	w_i = \mu_i + \sigma_i\,\varepsilon_i, \qquad \varepsilon_i\sim\cN(0,1),
\end{equation}
enabling low variance gradient estimates of~\eqref{eq:elbo} via
backpropagation.
A single noise draw is shared across all examples in a
mini batch, which inflates gradient variance at large batch sizes.

\subsubsection{Flipout}

\citet{wen2018flipout} reduce this variance by applying independent sign flip
matrices $\mathbf{r}_n,\mathbf{s}_n\in\{-1,+1\}^*$ to the shared
perturbation $\Delta\hat{\mathbf{W}}$ for each example $n$:
\begin{equation}
	\hat{\mathbf{W}}_n
	= \mathbf{W}_{\MAP} + \mathbf{s}_n \odot \Delta\hat{\mathbf{W}}
	\odot \mathbf{r}_n^{\!\top}.
\end{equation}

\subsubsection{MOPED Initialisation}

Training mean field BNNs from random initialisations is fragile.
\texttt{bayesian-torch}~\citep{intel2021bayesiantorch} addresses this via
MOPED: the posterior means are warm started from a
pre trained MAP checkpoint and the prior scale is set proportionally:
\begin{equation}
	\mu_i^{(0)} = \theta_{\MAP}^{(i)}, \qquad
	\sigma_i^{\text{prior}} = \delta\,|\theta_{\MAP}^{(i)}|, \quad \delta=0.5.
\end{equation}

\subsubsection{AvUC Loss}
\label{sec:theory:avuc}

Standard ELBO training minimises NLL but does not explicitly calibrate the
relationship between confidence and accuracy.
The
Accuracy vs Uncertainty Calibration loss, implemented in
\texttt{bayesian-torch}~\citep{intel2021bayesiantorch}, adds a direct
penalty for miscalibrated confidence.

\paragraph{AvUC objective.}
The AvUC loss aggregates the fractions of samples in the two
bad quadrants~\citep{intel2021bayesiantorch}:
\begin{equation}
	\mathcal{L}_{\mathrm{AvUC}}
	= \underbrace{\frac{n_{AU}}{n_{AU}+n_{AC}}}_{\text{over uncertain among correct}}
	+ \underbrace{\frac{n_{IC}}{n_{IC}+n_{IU}}}_{\text{overconfident among wrong}},
	\label{eq:avuc}
\end{equation}
where $n_X$ counts samples in quadrant $X$ within the current mini batch.
Each term lies in $[0,1]$; both are zero when calibration is perfect.

\subsection{Epistemic and Aleatoric Uncertainty Decomposition}
\label{sec:theory:unc_decomp}

For MC based BNNs, the total predictive uncertainty can be decomposed into
two orthogonal components~\citep{goan2020bayesian}:
\begin{align}
	\text{Total}     & = H\!\left[\frac{1}{S}\sum_{s=1}^{S} p(y|f_{\btheta_s}(\bm{x}))\right], \\
	\text{Aleatoric} & = \frac{1}{S}\sum_{s=1}^{S} H\!\left[p(y|f_{\btheta_s}(\bm{x}))\right], \\
	\text{Epistemic} & = \text{Total} - \text{Aleatoric}
	= I\!\left(y;\btheta\,|\,\bm{x},\cD\right),
	\label{eq:mi}
\end{align}

\subsection{Laplace Approximation}
\label{sec:theory:la}

The LA~\citep{daxberger2021laplace} builds a Gaussian posterior around
$\btheta_{\MAP}$ via second order Taylor expansion:
\begin{equation}
	p(\btheta|\cD) \approx \cN\!\bigl(\btheta;\,\btheta_{\MAP},\,\mathbf{H}^{-1}\bigr),
	\qquad
	\mathbf{H} = -\nabla^2\log p(\btheta_{\MAP}|\cD) \succ 0.
\end{equation}

\paragraph{Marginal likelihood tuning.}
The prior precision $\tau$ is selected by maximising the log evidence:
\begin{equation}
	\log p(\cD)
	\approx \log p(\cD|\btheta_{\MAP})
	+ \log p(\btheta_{\MAP})
	+ \tfrac{P}{2}\log 2\pi
	- \tfrac{1}{2}\log\det\mathbf{H},
	\label{eq:marglik}
\end{equation}
which is differentiable with respect to $\tau$ and requires no validation
data~\citep{daxberger2021laplace}. Grid search over
a discrete set of values on validation NLL is simpler but requires a
held out split.

\section{Evaluation Metrics}
\label{sec:metrics}

\paragraph{Brier Score.}
$\mathrm{BS}=\frac{1}{N}\sum_n\sum_c(\hat{p}_{nc}-\mathbf{1}[y_n=c])^2$;
bounded and less sensitive to extreme probabilities than NLL.

\paragraph{Mean Confidence / Mean Entropy / $1 - \text{Confidence}$.}
Mean max class probability, mean Shannon entropy $H$, and
$1-\max_c\hat{p}_c\in[0,0.5]$.

\paragraph{Uncertainty calibration U MAE, U MSE, U Corr, Max Gap.}
Compares expected uncertainty against empirical error rate
per bin. Pearson correlation and max bin gap assess worst case alignment.

\paragraph{Epistemic uncertainty OOD maps.}
Mutual information $I(y;\btheta|\bm{x},\cD)$ estimated via
equation~\eqref{eq:mi} with $S=50$ MC samples.

\section{Shared Infrastructure}
\label{sec:shared}

Both notebooks share a \texttt{shared} utility package. The MAP baseline uses Adam with early stopping.

\section{Two Moons Classification}
\label{sec:twomoons}

\subsection{Dataset and OOD Probe Design}

2000 2D points from two interleaved half crescents.
To probe out of distribution behaviour, 200 probe points are
sampled in four corner clusters deliberately outside the training manifold.

\subsection{\texttt{bayesian-torch}: Variational BNN}
\label{sec:bnn}

\begin{lstlisting}[caption={BNN conversion and ELBO AvUC training.}]
dnn_to_bnn(net, {"prior_mu": 0.0, "prior_sigma": 1.0,
    "posterior_mu_init": 0.0, "posterior_rho_init": -3.0,
    "type": "Reparameterization",
    "moped_enable": True, "moped_delta": 0.5})

loss = CE(mean_logits, labels) + get_kl_loss(net) / N + avu
\end{lstlisting}

Table~\ref{tab:bnn_main_subset} reports the relevant calibration statistics for all \texttt{bayesian-torch} configurations.

\begin{table}[H]
	\centering
	\caption{Relevant calibration metrics for bayesian torch variants on Two Moons.}
	\label{tab:bnn_main_subset}
	\small\setlength{\tabcolsep}{4pt}
	\begin{tabular}{llccccccc}
		\toprule
		Type    & Init & Loss & BS    & $1 - c$ & U MAE & U MSE & U Corr & MaxGap \\
		\midrule
		Flipout & scr  & ELBO & 0.058 & 0.094   & 0.018 & 0.458 & 0.988  & 0.069  \\
		Flipout & scr  & AvUC & 0.059 & 0.072   & 0.014 & 0.840 & 0.953  & 0.134  \\
		Flipout & MOP  & ELBO & 0.059 & 0.089   & 0.019 & 0.882 & 0.935  & 0.113  \\
		Flipout & MOP  & AvUC & 0.059 & 0.072   & 0.010 & 0.294 & 0.990  & 0.078  \\
		BBB     & scr  & ELBO & 0.058 & 0.088   & 0.019 & 0.582 & 0.983  & 0.051  \\
		BBB     & scr  & AvUC & 0.095 & 0.153   & 0.028 & 0.918 & 0.965  & 0.102  \\
		BBB     & MOP  & ELBO & 0.058 & 0.087   & 0.015 & 0.556 & 0.967  & 0.130  \\
		BBB     & MOP  & AvUC & 0.057 & 0.067   & 0.014 & 0.900 & 0.941  & 0.123  \\
		\bottomrule
	\end{tabular}
\end{table}

\begin{table}[H]
	\centering
	\caption{OOD epistemic uncertainty decomposition for all bayesian torch variants.}
	\label{tab:ood}
	\small\setlength{\tabcolsep}{4pt}
	\begin{tabular}{lccccccc}
		\toprule
		Variant       & ID $1 - c$ & OOD $1 - c$ & ID Alea & OOD Alea & ID Epist & OOD Epist & $\Delta$Epist \\
		\midrule
		Flip scr ELBO & 0.094      & 0.005       & 0.238   & 0.028    & 0.003    & 0.001     & $-0.003$      \\
		Flip scr AvUC & 0.072      & 0.003       & 0.177   & 0.012    & 0.008    & 0.001     & $-0.007$      \\
		Flip MOP ELBO & 0.089      & 0.010       & 0.213   & 0.034    & 0.009    & 0.003     & $-0.006$      \\
		Flip MOP AvUC & 0.073      & 0.004       & 0.150   & 0.018    & 0.036    & 0.004     & $-0.032$      \\
		BBB scr ELBO  & 0.089      & 0.003       & 0.227   & 0.018    & 0.002    & 0.000     & $-0.002$      \\
		BBB scr AvUC  & 0.153      & 0.032       & 0.359   & 0.123    & 0.003    & 0.002     & $-0.001$      \\
		BBB MOP ELBO  & 0.086      & 0.011       & 0.202   & 0.039    & 0.014    & 0.004     & $-0.009$      \\
		BBB MOP AvUC  & 0.067      & 0.007       & 0.140   & 0.027    & 0.035    & 0.009     & $-0.026$      \\
		\bottomrule
	\end{tabular}
\end{table}

Please refer to the Appendix for histograms, boundary visualizations, and full tables.

\subsection{\texttt{laplace-torch}: Laplace Approximation}
\label{sec:la}

\begin{lstlisting}[caption={Laplace post hoc workflow.}]
la = Laplace(model, 'classification',
             subset_of_weights='last_layer', hessian_structure='kron')
la.fit(train_loader)
la.optimize_prior_precision(method='gridsearch', val_loader=val_loader,
                             pred_type='glm', link_approx='probit')
probs = la(X_test, pred_type='glm', link_approx='probit')
\end{lstlisting}

Table~\ref{tab:la_main_subset} reports the relevant subset of metrics for the requested Laplace configurations including Diagonal all weights, Full Last Layer, and Kronecker Last Layer under both Marginal Likelihood and Gridsearch tuning.

\begin{table}[H]
	\centering
	\caption{Relevant calibration metrics for laplace torch on Two Moons.}
	\label{tab:la_main_subset}
	\small\setlength{\tabcolsep}{4pt}
	\begin{tabular}{llccccccc}
		\toprule
		Variant     & Scope & PP     & BS    & $1 - c$ & U MAE & U MSE & U Corr & MaxGap \\
		\midrule
		Diag all GS & all   & $10^4$ & 0.056 & 0.079   & 0.010 & 0.806 & 0.959  & 0.197  \\
		Full LL GS  & LL    & $10^4$ & 0.056 & 0.079   & 0.010 & 0.806 & 0.959  & 0.197  \\
		Kron LL GS  & LL    & $10^4$ & 0.056 & 0.079   & 0.010 & 0.806 & 0.959  & 0.197  \\
		Diag all ML & all   & 8.73   & 0.057 & 0.088   & 0.016 & 0.817 & 0.955  & 0.196  \\
		Full LL ML  & LL    & 1.79   & 0.062 & 0.145   & 0.070 & 6.49  & 0.935  & 0.136  \\
		Kron LL ML  & LL    & 2.04   & 0.061 & 0.138   & 0.064 & 5.50  & 0.938  & 0.141  \\
		\bottomrule
	\end{tabular}
\end{table}

Please refer to the Appendix for epistemic maps, full metrics, and visualizations.

\section{Discussion}
\label{sec:discussion}

MOPED is a geometric anchor. The BBB scratch AvUC failure demonstrates that AvUC calibration alone is not sufficient.
Under marglik, full matrix GGN curvature is poorly conditioned and selects pathological prior precisions.
All weights LA marglik variants retain some OOD epistemic sensitivity, while last layer LA and all SVI variants collapse to near zero outside the boundary.

\appendix
\section{Appendix: Full Metric Tables and Visualizations}
\label{sec:app}

\begin{sidewaystable}
	\centering
	\caption{All calibration metrics for bayesian torch variants on Two Moons.}
	\label{tab:bnn_main_full}
	\scriptsize\setlength{\tabcolsep}{2.5pt}
	\begin{tabular}{
			L{1.5cm} C{0.6cm} C{0.65cm} C{0.65cm} C{0.7cm} C{0.7cm} C{0.7cm}
			C{0.7cm} C{0.7cm} C{0.65cm} C{0.7cm} C{0.95cm} C{0.7cm}
			C{0.75cm} C{0.7cm} C{0.75cm} C{0.5cm} C{0.9cm} C{0.85cm}
		}
		\toprule
		Type    & Init  & Loss  & Acc    & NLL    & BS    & ECE     & Conf  & Ent   & $1{-}c$
		        & U MAE & U MSE & U Corr & MaxGap & ClECE & MActErr & Ep    & Fit   & Inf                                                                         \\
		\midrule
		Flipout & scr   & ELBO  & 0.921  & 0.198  & 0.058 & 0.018   & 0.906 & 0.241 & 0.094   & 0.018 & 0.458 & 0.988 & 0.069 & 0.018 & 0.080 & 39 & 19.1 & 0.064 \\
		Flipout & scr   & AvUC  & 0.918  & 0.198  & 0.059 & 0.014   & 0.928 & 0.184 & 0.072   & 0.014 & 0.840 & 0.953 & 0.134 & 0.014 & 0.082 & 38 & 80.4 & 0.065 \\
		Flipout & MOP   & ELBO  & 0.918  & 0.199  & 0.059 & 0.019   & 0.911 & 0.222 & 0.089   & 0.019 & 0.882 & 0.935 & 0.113 & 0.019 & 0.083 & 39 & 20.0 & 0.063 \\
		Flipout & MOP   & AvUC  & 0.920  & 0.201  & 0.059 & 0.010   & 0.928 & 0.184 & 0.072   & 0.010 & 0.294 & 0.990 & 0.078 & 0.010 & 0.081 & 40 & 90.1 & 0.064 \\
		BBB     & scr   & ELBO  & 0.923  & 0.197  & 0.058 & 0.019   & 0.911 & 0.229 & 0.088   & 0.019 & 0.582 & 0.983 & 0.051 & 0.019 & 0.078 & 50 & 23.2 & 0.051 \\
		BBB     & scr   & AvUC  & 0.866  & 0.311  & 0.095 & 0.028   & 0.847 & 0.362 & 0.153   & 0.028 & 0.918 & 0.965 & 0.102 & 0.028 & 0.134 & 32 & 73.3 & 0.050 \\
		BBB     & MOP   & ELBO  & 0.920  & 0.195  & 0.058 & 0.015   & 0.913 & 0.217 & 0.087   & 0.015 & 0.556 & 0.967 & 0.130 & 0.015 & 0.081 & 36 & 15.9 & 0.055 \\
		BBB     & MOP   & AvUC  & 0.923  & 0.192  & 0.057 & 0.014   & 0.933 & 0.173 & 0.067   & 0.014 & 0.900 & 0.941 & 0.123 & 0.014 & 0.078 & 33 & 68.7 & 0.051 \\
		\bottomrule
	\end{tabular}
\end{sidewaystable}

\begin{sidewaystable}
	\centering
	\caption{All calibration metrics for laplace torch variants on Two Moons.}
	\label{tab:la_main_full}
	\scriptsize\setlength{\tabcolsep}{2.5pt}
	\begin{tabular}{
			L{1.9cm} C{0.6cm} C{0.55cm} C{0.7cm} C{0.7cm} C{0.7cm} C{0.7cm}
			C{0.7cm} C{0.7cm} C{0.65cm} C{0.7cm} C{0.95cm} C{0.7cm}
			C{0.75cm} C{0.7cm} C{0.95cm} C{0.75cm} C{0.95cm}
		}
		\toprule
		Variant     & Hess  & Sc    & PP     & NLL    & BS    & ECE   & Conf  & Ent   & $1{-}c$
		            & U MAE & U MSE & U Corr & MaxGap & ClECE & Fit   & Tune  & Inf                                                                           \\
		\midrule
		Diag all GS & diag  & all   & $10^4$ & 0.190  & 0.056 & 0.010 & 0.921 & 0.194 & 0.079   & 0.010 & 0.806 & 0.959 & 0.197 & 0.010 & 0.23 & 4.01 & 0.021 \\
		Full all GS & full  & all   & $10^4$ & 0.190  & 0.056 & 0.010 & 0.921 & 0.194 & 0.079   & 0.010 & 0.806 & 0.959 & 0.197 & 0.010 & 1263 & 11.4 & 0.034 \\
		Kron all GS & kron  & all   & $10^4$ & 0.190  & 0.056 & 0.010 & 0.921 & 0.194 & 0.079   & 0.010 & 0.806 & 0.959 & 0.197 & 0.010 & 0.54 & 5.29 & 0.022 \\
		Full LL GS  & full  & LL    & $10^4$ & 0.190  & 0.056 & 0.010 & 0.921 & 0.194 & 0.079   & 0.010 & 0.806 & 0.959 & 0.197 & 0.010 & 7.12 & 1.52 & 0.001 \\
		Kron LL GS  & kron  & LL    & $10^4$ & 0.190  & 0.056 & 0.010 & 0.921 & 0.194 & 0.079   & 0.010 & 0.806 & 0.959 & 0.197 & 0.010 & 0.43 & 1.85 & 0.001 \\
		Diag all ML & diag  & all   & 8.73   & 0.192  & 0.057 & 0.016 & 0.912 & 0.219 & 0.088   & 0.016 & 0.817 & 0.955 & 0.196 & 0.016 & 0.20 & 0.06 & 0.021 \\
		Full all ML & full  & all   & 0.52   & 0.307  & 0.082 & 0.152 & 0.771 & 0.500 & 0.229   & 0.152 & 26.0  & 0.913 & 0.218 & 0.152 & 1122 & 1.88 & 0.034 \\
		Kron all ML & kron  & all   & 1.42   & 0.238  & 0.065 & 0.086 & 0.841 & 0.377 & 0.159   & 0.086 & 9.07  & 0.923 & 0.152 & 0.086 & 0.53 & 0.10 & 0.022 \\
		Full LL ML  & full  & LL    & 1.79   & 0.226  & 0.062 & 0.070 & 0.855 & 0.348 & 0.145   & 0.070 & 6.49  & 0.935 & 0.136 & 0.070 & 7.13 & 0.09 & 0.001 \\
		Kron LL ML  & kron  & LL    & 2.04   & 0.221  & 0.061 & 0.064 & 0.862 & 0.334 & 0.138   & 0.064 & 5.50  & 0.938 & 0.141 & 0.064 & 0.44 & 0.08 & 0.001 \\
		\bottomrule
	\end{tabular}
\end{sidewaystable}

\begin{figure}[H]
	\centering
	\includegraphics[width=0.60\linewidth]{ood_points.png}
	\caption{Training data and OOD probe points.}
	\label{fig:ood_points}
\end{figure}

\begin{figure}[H]
	\centering
	\begin{subfigure}{0.48\linewidth}
		\includegraphics[width=\linewidth]{bbb_moped_avuc_boundary.png}
		\caption{BBB MOPED AvUC boundary}
	\end{subfigure}\hfill
	\begin{subfigure}{0.48\linewidth}
		\includegraphics[width=\linewidth]{bbb_scratch_avuc_boundary.png}
		\caption{BBB scratch AvUC boundary}
	\end{subfigure}
	\\[6pt]
	\begin{subfigure}{0.48\linewidth}
		\includegraphics[width=\linewidth]{bbb_moped_avuc_uncertainty.png}
		\caption{BBB MOPED AvUC uncertainty map}
	\end{subfigure}\hfill
	\begin{subfigure}{0.48\linewidth}
		\includegraphics[width=\linewidth]{bbb_scratch_avuc_uncertainty.png}
		\caption{BBB scratch AvUC uncertainty}
	\end{subfigure}
	\caption{MOPED vs scratch initialisation with AvUC loss.}
	\label{fig:bnn_boundaries}
\end{figure}

\begin{figure}[H]
	\centering
	\begin{subfigure}{0.48\linewidth}
		\includegraphics[width=\linewidth]{bbb_moped_avuc_histogram.png}
		\caption{BBB MOPED AvUC histogram}
	\end{subfigure}\hfill
	\begin{subfigure}{0.48\linewidth}
		\includegraphics[width=\linewidth]{bbb_moped_avuc_reliability.png}
		\caption{BBB MOPED AvUC reliability}
	\end{subfigure}
	\\[6pt]
	\begin{subfigure}{0.48\linewidth}
		\includegraphics[width=\linewidth]{bbb_scratch_avuc_histogram.png}
		\caption{BBB scratch AvUC histogram}
	\end{subfigure}\hfill
	\begin{subfigure}{0.48\linewidth}
		\includegraphics[width=\linewidth]{bbb_moped_elbo_reliability.png}
		\caption{BBB MOPED ELBO reliability}
	\end{subfigure}
	\caption{Probability histograms and reliability diagrams.}
	\label{fig:bnn_histrel}
\end{figure}

\begin{figure}[H]
	\centering
	\begin{subfigure}{0.48\linewidth}
		\includegraphics[width=\linewidth]{bbb_moped_avuc_ood_epistemic.png}
		\caption{BBB MOPED AvUC}
	\end{subfigure}\hfill
	\begin{subfigure}{0.48\linewidth}
		\includegraphics[width=\linewidth]{bbb_scratch_elbo_ood_epistemic.png}
		\caption{BBB scratch ELBO}
	\end{subfigure}
	\caption{OOD epistemic uncertainty.}
	\label{fig:bnn_ood}
\end{figure}

\begin{figure}[H]
	\centering
	\begin{subfigure}{0.48\linewidth}
		\includegraphics[width=\linewidth]{marglik_full_all_weights_boundary.png}
		\caption{Full all marglik boundary}
	\end{subfigure}\hfill
	\begin{subfigure}{0.48\linewidth}
		\includegraphics[width=\linewidth]{marglik_full_all_weights_reliability.png}
		\caption{Full all marglik reliability}
	\end{subfigure}
	\caption{Marginal likelihood failure for Full all LA.}
	\label{fig:la_failure}
\end{figure}

\begin{figure}[H]
	\centering
	\begin{subfigure}{0.48\linewidth}
		\includegraphics[width=\linewidth]{marglik_kron_last_layer_uncertainty.png}
		\caption{Kron LL marglik}
	\end{subfigure}\hfill
	\begin{subfigure}{0.48\linewidth}
		\includegraphics[width=\linewidth]{marglik_kron_all_weights_ood_epistemic.png}
		\caption{Kron all marglik OOD epistemic}
	\end{subfigure}
	\caption{Uncertainty patterns under marginal likelihood tuning.}
	\label{fig:la_stripe}
\end{figure}

% conference papers do not normally have an appendix


% use section* for acknowledgement
% \section*{Acknowledgment}
% The authors would like to thank...




% trigger a \newpage just before the given reference
% number - used to balance the columns on the last page
% adjust value as needed - may need to be readjusted if
% the document is modified later
%\IEEEtriggeratref{8}
% The "triggered" command can be changed if desired:
%\IEEEtriggercmd{\enlargethispage{-5in}}



% references section
\bibliographystyle{plainnat}
\bibliography{bib.bib}


% that's all folks
\end{document}
